% ============================================================

% The regulatory timeline lives in the appendix (app:timeline); this frame
% reaches it via the Timeline button.
\begin{frame}{The Judicialization of Electoral Competition}
\hypertarget{src:motivation}{}
\small
\begin{itemize}\itemsep1pt
  \item Electoral governance operates on three levels: rulemaking, rule application, and adjudication \ns{\citep{mozaffar2002}}. In Brazil one branch holds all three. The \emph{Justi\c{c}a Eleitoral} fixes terms of the contest by ruling (councillor seats, party loyalty, coalition verticalization), then resolves the disputes those terms generate \ns{\citep{marchetti2012,marchetticortez2009}}.\hfill\hyperlink{app:theory}{\beamergotobutton{Theory}}

  \item The arena is not merely obeyed; it is used. A third of mayoral candidates are party to an electoral lawsuit, and a ninth of campaign money goes to lawyers. Most suits are filed in the final weeks, disproportionately by the runner-up against leading candidates; suits seeking disqualification almost never succeed, and candidates feature them in their media campaigns \ns{\citep{chin2025}}. Attacks target candidates with an electoral advantage \ns{\citep{nakaguma2025}}.

  \item Voters respond to what the courts say: a conviction disclosed before the vote costs some thirteen points of vote share \ns{\citep{assumpcao2019}}. Both literatures are candidate-level; what litigation does to the race itself is not measured.

  \item Most of the \TotalAllLawsuits{} million filings across 2020 and 2024 is routine paperwork; the \textcolor{myblue}{adversarial core} fought between political actors stays flat and recomposes as new fronts open. The object is \textcolor{myblue}{intensity}, not a trend.\hfill\hyperlink{app:recomp}{\beamergotobutton{Recomposition}}\,\hyperlink{app:timeline}{\beamergotobutton{Timeline}}
\end{itemize}
\vspace{-2pt}
\begin{block}{Research question}
\footnotesize As a race faces more adversarial litigation, what happens to the local election? Do voters change how they participate and how they mark the ballot? Does the litigated campaign fall differently on different candidates? Judicialization reaches the election as noise around the contest, and so through how the electorate reads it.
\end{block}
\end{frame}
